\section*{Physique (13 points)}
\section*{Exercice 1 (3,5 points): Propagation des ondes}
L'étude de la propagation des ondes mécaniques et des ondes lumineuses permet de déterminer certaines caractéristiques des ondes et propriétés des milieux de propagation.\\
Cet exercice vise l'étude de la propagation d'une onde sonore dans l'air et l'étude de la dispersion de la lumière.

\section*{1. Détermination de la vitesse de propagation d'une onde sonore}
Un haut-parleur relié à un générateur basse fréquence ( $G B F$ ) émet un signal sonore de fréquence $N$. Ce signal est capté par un microphone situé le long de l'axe ( $O x$ ). Ce microphone est relié à un oscilloscope (figure 1).\\
La figure (2) donne l'enregistrement de deux signaux captés par le microphone pour deux positions successives $x_{1}$ et $x_{2}$.\\
Le signal (a) correspond à $x_{1}=20 \mathrm{~cm}$. Le signal (b) correspond à $x_{2}=36,7 \mathrm{~cm}$, et apparait pour la première fois en phase avec le signal (a).

Figure 1

\begin{figure}[h]
\begin{center}
  \includegraphics[alt={},max width=\textwidth]{3722cc87-920c-4836-a0f1-9718f4bdb853-3_570_807_1809_1055}
\captionsetup{labelformat=empty}
\caption{Figure 2}
\end{center}
\end{figure}

Donnée : Sensibilité horizontale : $0,2 \mathrm{~ms} . \mathrm{div}^{-1}$\\
1.1. Déterminer la valeur de la fréquence $N$.\\
1.2. Déterminer la valeur de longueur d'onde $\lambda$ de l'onde sonore.\\
1.3. Déterminer la valeur de la vitesse v de propagation de cette onde.

\section*{2. Identification d'un milieu dispersif}
Le tableau ci-dessous donne les longueurs d'onde dans le vide de deux radiations monochromatiques (violette et bleue), et les indices de réfraction correspondants à chaque longueur d'onde pour trois milieux de propagation : l'air, le verre crown et le verre flint.

\begin{center}
\begin{tabular}{|l|c|c|}
\hline
Couleur de la radiation & violette & bleue \\
\hline
Longueur d'onde dans le vide & $\lambda_{0 v}=486,1 \mathrm{~nm}$ & $\lambda_{0 b}=589 \mathrm{~nm}$ \\
\hline
Indice de réfraction de l'air & $n_{a}=1$ & $n_{a}=1$ \\
\hline
Indice de réfraction du verre crown & $n_{c}=1,522$ & $n_{c}=1,517$ \\
\hline
Indice de réfraction du verre flint & $n_{f}=1,682$ & $n_{f}=1,666$ \\
\hline
\end{tabular}
\end{center}

Donnée : $c=3.10^{8} \mathrm{~m} . \mathrm{s}^{-1}$\\
2.1. Déterminer la valeur de la fréquence de la radiation bleue.\\
2.2. Établir la relation entre l'indice $n$ d'un milieu, la longueur d'onde $\lambda$, la fréquence $v$ d'une radiation et la célérité $c$ de la lumière dans le vide.\\
2.3. Parmi les trois milieux proposés, indiquer en justifiant, ceux qui sont dispersifs.\\
2.4. Déterminer la valeur de la longueur d'onde $\lambda_{b}$ de la radiation bleue dans le verre flint.

\section*{Exercice 2 (5,5 points): Comportement d'un condensateur dans un circuit électrique}
Les associations de composants électriques comme la bobine, le condensateur et le conducteur ohmique conduisent à différents dipôles électriques comme $R C$ et $R L C$ qui, placés dans des circuits engendrent des phénomènes tel que la charge ou la décharge du condensateur et les oscillations électriques libres...\\
Cet exercice vise :

\begin{itemize}
  \item l'étude de la réponse d'un dipôle $R C$ à un échelon de tension ;
  \item l'étude énergétique d'un circuit oscillant $L C$.
\end{itemize}

On étudie le comportement d'un condensateur dans deux situations (a) et (b) différentes en utilisant le montage de la figure (1) qui comporte :

\begin{itemize}
  \item un générateur idéal de tension de f.e.m $E$;
  \item un condensateur de capacité C;
  \item deux conducteurs ohmiques de résistances $R$ et $R^{\prime}$;
  \item une bobine d'inductance $L$ et de résistance négligeable;
  \item un interrupteur $K$ à double position.
\end{itemize}

Donnée : $R=100 \Omega$

\begin{figure}[h]
\begin{center}
  \includegraphics[alt={},max width=\textwidth]{3722cc87-920c-4836-a0f1-9718f4bdb853-4_424_577_1731_1315}
\captionsetup{labelformat=empty}
\caption{Figure 1}
\end{center}
\end{figure}

\section*{Partie 1 : Étude du comportement du condensateur dans la situation (a)}
À l'instant $t_{0}=0$, on place l'interrupteur $K$ dans la position (1).

\begin{enumerate}
  \item Quel est l'intérêt du montage dans ce cas?
  \item En utilisant la loi d'additivité des tensions, montrer que l'intensité $i(t)$ du courant qui circule dans le circuit est liée à la charge $q(t)$ du condensateur par la relation : $i=-\frac{1}{R \cdot C} \cdot q+\frac{E}{R}$.
  \item À l'aide d'un système d'acquisition convenable, on obtient le graphe de la figure (2) qui représente l'évolution de $i$ en fonction de $q$.\\
En exploitant le graphe, déterminer les valeurs de :\\
a. l'intensité maximale $I_{0}$ du courant électrique;\\
b. la f.e.m $E$;\\
c. la constante de temps $\tau$ du circuit ;\\
d. la charge maximale $Q_{\text {max }}$ du condensateur à la fin de sa charge.
\end{enumerate}

\begin{figure}[h]
\begin{center}
  \includegraphics[alt={},max width=\textwidth]{3722cc87-920c-4836-a0f1-9718f4bdb853-5_511_666_310_1244}
\captionsetup{labelformat=empty}
\caption{Figure 2}
\end{center}
\end{figure}

\section*{Partie 2 : Étude du comportement du condensateur dans la situation (b)}
Jne fois le condensateur totalement chargé sous la tension $u_{C 0}=E$ dans la situation (a), on bascule 'interrupteur $K$ en position (2) à l'instant $t_{0}=0$.\\
-a figure (3) donne les variations de la tension $u_{C}(t)$ aux bornes du condensateur.\\
l. Nommer le régime d'oscillations mis en évidence par le graphe de la figure (3).\\
-. Expliquer du point de vue énergétique le régime l'oscillations dans le circuit.\\
3. On note respectivement $\mathscr{E}$ et $\mathscr{E}$ les énergies otales du circuit aux instants $t_{0}=0$ et $t_{1}=188 \mathrm{~ms}$. La variation de l'énergie totale du circuit entre $t_{0}$ it $t_{1}$ est $\Delta \mathscr{E}=-10,5.10^{-4} J$.\\
3.1. On désigne par $u_{C 1}$ la tension aux bornes du ondensateur à l'instant $t_{1}$.\\
Montrer que la capacité du condensateur peut

\begin{figure}[h]
\begin{center}
  \includegraphics[alt={},max width=\textwidth]{3722cc87-920c-4836-a0f1-9718f4bdb853-5_497_835_1242_1114}
\captionsetup{labelformat=empty}
\caption{Figure 3}
\end{center}
\end{figure}

Galculer la valeur de $C$.\\
3.2. Le condensateur utilisé peut être remplacé par deux condensateurs identiques ayant chacun la zapacité $C_{0}$ montés en parallèle. Déterminer la valeur de $C_{0}$.\\
3.3. On suppose que la pseudo-période $T$ est égale à la période propre $T_{0}$ des oscillations libres non imorties. Déterminer la valeur de $L$ (On prend $\pi^{2}=10$ ).

\section*{Exercice 3 (4 points) : Mouvement d'un solide sur un plan incliné}
Les mouvements des systèmes mécaniques sont généralement régis par les lois de Newton. L'état de mouvement de ces systèmes dépend des actions mécaniques exercées et des conditions initiales.\\
Cet exercice vise la détermination de certaines grandeurs lors du mouvement d'un solide sur un plan incliné.

On considère un solide ( $S$ ) de masse $m$, susceptible de glisser selon la ligne de plus grande pente d'un plan incliné faisant un angle $\alpha$ avec l'horizontal.\\
Le solide ( $S$ ) démarre avec une vitesse initiale $\overrightarrow{\mathrm{v}}_{0}$, à l'instant $t_{0}=0$ à partir de la position $O$. Au cours de son mouvement le long du trajet $O A$, le solide est soumis à des frottements modélisés par\\
une force $\vec{f}$ constante de même direction que le vecteur vitesse et de sens opposé.\\
On étudie le mouvement du centre d'inertie $G$ du solide ( $S$ ) dans le repère ( $O, \vec{i}$ ) lié à la Terre supposé galiléen (figure 1).\\
L'abscisse de $G$ à $t_{0}=0$ est $x_{G}=x_{0}=0$.\\
Données : $m=500 \mathrm{~g} \quad ; \quad g=10 \mathrm{~m} . \mathrm{s}^{-2} \quad ; \quad \alpha=20^{\circ}$

\begin{enumerate}
  \item En appliquant la deuxième loi de Newton, montrer que l'équation différentielle vérifiée par $x_{G}$ s'écrit : $\frac{d^{2} x_{G}}{d t^{2}}=g . \sin \alpha-\frac{f}{m}$.
  \item La figure (2) donne l'évolution de la vitesse $\mathrm{v}(t)$ de $G$.
\end{enumerate}

\begin{figure}[h]
\begin{center}
  \includegraphics[alt={},max width=\textwidth]{3722cc87-920c-4836-a0f1-9718f4bdb853-6_305_414_411_1503}
\captionsetup{labelformat=empty}
\caption{Figure 1}
\end{center}
\end{figure}

2.1. Déterminer graphiquement les valeurs de l'accélération $a_{G}$ et de la vitesse initiale $\mathrm{v}_{0}$ du mouvement de $G$.\\
2.2. Écrire l'équation horaire $x(t)$ du mouvement de $G$.\\
2.3. Calculer l'intensité de la force $\vec{f}$.\\
3. Après passage du solide ( $S$ ) par la position $A$ avec la vitesse $\mathrm{v}_{A}=6 \mathrm{~m} \cdot \mathrm{~s}^{-1}$, celui-ci n'est plus soumis à la force de frottement $\vec{f}$, il passe ensuite par une position $B$ après avoir parcouru une distance $A B$.\\
3.1. Déterminer la nature du mouvement de $G$ après son passage par $A$.\\
3.2. On choisit la position $A$ comme nouvelle origine des

\begin{figure}[h]
\begin{center}
  \includegraphics[alt={},max width=\textwidth]{3722cc87-920c-4836-a0f1-9718f4bdb853-6_485_643_918_1283}
\captionsetup{labelformat=empty}
\caption{Figure 2}
\end{center}
\end{figure}

Le centre d'inertie $G$ du solide ( $S$ ) passe par la position $B$ avec une vitesse $\overrightarrow{\mathrm{V}}_{B}$ à l'instant $t=1 \mathrm{~s}$.\\
Déterminer :\\
a. la valeur de la distance $A B$.\\
b. la valeur de la vitesse $\mathrm{v}_{B}$.\\
3.3. Déterminer l'intensité de la force $\vec{R}$ exercée par le plan incliné sur le solide ( $S$ ).


